
\paragraph{另一种违约解释的对照}
若原计划费照付，再对退购加收50\%违约费，则应改为
\begin{equation}
 \phi_{+}(q,r)=q+1.5(r-q)^++0.5(q-r)^+,
 \qquad \phi_{+}(q,r)-\phi(q,r)=(q-r)^+.
\end{equation}
以下将同一334日主轨迹按两种规则重计费，以隔离口径变化；它不是新口径下重新优化的年度最优费用。
\begin{table}[htbp]
\centering
\caption{同终点原策略在两种退购计费解释下的全年账单；新口径未重新优化}
\begin{tabular}{lrrr}
\toprule
问题 & 原口径（元） & 附加违约费口径（元） & 增幅（\%）\\
\midrule
第二问 & 13844170.06 & 13844170.06 & 0.000\\
第三问 & 13461842.35 & 14073209.01 & 4.541\\
第四问日前 & 14654611.38 & 14654611.38 & 0.000\\
第四问日内 & 14252083.03 & 14907515.48 & 4.599\\
\bottomrule
\end{tabular}
\end{table}

在非负价格、免费且不受限的弃电条件下，将$r<q$改为$r'=q$、保持原电池动作并增加弃电$q-r$，物理平衡不变，新口径费用减少$0.5p(q-r)$。因此该条件下总存在不退购的最优动作，结论不覆盖负价或弃电受限情形。

另在1月25--31日保持初末6000 kWh、预测和权限一致，改变上层计费上图并重新滚动执行：问题三新口径费用345349.60元，较原策略直接重计费的350544.87元少5195.27元；问题四日内为384732.76元，较390984.78元少6252.02元。新策略的退购量均为数值零。这是开发窗口的计费敏感性，不替换主口径全年结果。
